\subsection{Понятие и особенности ончейн-данных}

Под ончейн-данными в настоящей работе понимаются данные, зафиксированные в
блоках канонической цепи блокчейна, а также связанные с ними производные
структуры, необходимые для прикладной интерпретации состояния сети. К таким
данным относятся сведения о блоках, транзакциях, входах и выходах транзакций,
адресах, текущем наборе \gls{utxo}, исторических балансах и состоянии
неподтвержденных транзакций в mempool.

Существенная особенность ончейн-данных заключается в их событийной природе.
Состояние адресов и активов не хранится в блокчейне в виде готовых агрегатов,
а восстанавливается на основе последовательной обработки транзакций и анализа
связей между их входами и выходами. Это означает, что для прикладного
использования данных недостаточно иметь доступ только к сырому журналу блоков:
требуется дополнительный слой нормализации, агрегации и индексирования.

Другой особенностью является зависимость интерпретации данных от состояния
канонической цепи. При возникновении reorg часть ранее подтвержденных блоков и
транзакций может потерять канонический статус, что приводит к необходимости
повторной актуализации адресных индексов, агрегированных балансов и набора
непотраченных выходов. Следовательно, система работы с ончейн-данными должна
поддерживать не только накопление данных, но и их согласованную коррекцию.

Для прикладных сервисов особое значение имеют следующие свойства
ончейн-данных:
\begin{enumerate}
    \item большой объем и непрерывное пополнение по мере роста цепи;
    \item необходимость восстановления производных представлений из первичных записей;
    \item чувствительность к изменению канонической ветки;
    \item различие между подтвержденными данными и временным состоянием mempool;
    \item высокая стоимость повторных выборок при отсутствии специализированных индексов.
\end{enumerate}

Таким образом, ончейн-данные представляют собой не просто набор записей
блокчейна, а сложный предметный массив, требующий специализированных средств
структурирования, хранения и выдачи.
